\section{DGLAP and running-coupling terms against the literature}
\label{sec:dglap-lit}

\noindent\textbf{Scope.} The DGLAP term and the running-coupling (RC) term of the final result, and the contribution they come from (Group~II Row~IV). Literature used, all read in this session from the arXiv sources: Chirilli, Xiao, Yuan (CXY) arXiv:1203.6139 and 1112.1061; Altinoluk, Armesto, Beuf, Kovner, Lublinsky (AABKL) arXiv:1411.2869; Kovner, Lublinsky, Skokov, Zhao (KLSZ) arXiv:2308.15545; Horowitz, Kovchegov (HK) arXiv:1009.0545; Kovchegov, Weigert (KW) hep-ph/0612071 and arXiv:0712.3732. Equations of these papers are cited by their \LaTeX{} labels, because the equation numbers of the printed versions were not available. Scripts: \texttt{finalchk/dglap2/verify.py}, \texttt{checks.py}. Notation: $\mathcal{C}\equiv\frac{1}{(2\pi)^{3}}\frac{g^{2}}{2\pi^{2}}\frac{g^{2}}{8\pi^{2}}\frac{1}{k^{+}}$, and
\begin{align}
\mathcal{F}(\xi)&\equiv\frac{1}{\xi^{2}}\int_{x,y,x',y'}e^{-ik\cdot(y-y')\frac{1}{\xi}}\Big[2\gamma_{E}+\log\frac{(y'-y)^{2}\mu^{2}}{4\xi^{2}}\Big]\mathcal{K}(x-y)\cdot\mathcal{K}(x'-y')\,[U(y')-U(x')]^{ce'}[U(y)-U(x)]^{ce}\rho^{e'}(x')\rho^{e}(x),\notag\\
\mathcal{G}_{0}(\xi)&\equiv\text{the same without the bracket},\qquad \mathcal{X}\equiv\int_{x,x',z}\mathcal{K}(x-z)\cdot\mathcal{K}(x'-z)\,[U(z)-U(x')]^{ce'}[U(z)-U(x)]^{ce}\rho^{e'}(x')\rho^{e}(x),
\end{align}
so that $N_{\rm DGLAP}=-\mathcal{C}\int_{0}^{1}d\xi\,\widetilde{\mathcal P}_{gg}(\xi)\mathcal{F}(\xi)$, $N_{\rm RC}=+\mathcal{C}\frac{\beta_{0}}{2}\mathcal{F}(1)$, $\mathcal{F}(0)=-\frac{4\pi}{k^{2}}\mathcal{X}$, $\mathcal{G}_{0}(0)=0$ for $k\neq0$. Here $\xi=k^{+}/(p^{+}+k^{+})$ is the momentum fraction of the measured gluon in its parent; the parent sits at $y=z+\xi(w-z)$ ($w$: measured gluon, $z$: integrated gluon), $U(z)$ cancels by unitarity, and the parent carries the Wilson-line structure of the LO amplitude. Row~IV is therefore a \emph{final-state} splitting $g\to gg$ of a gluon produced at LO.

\subsection{The splitting function}

\noindent\textbf{Known result} (CXY 1203.6139, Sec.\ ``The gluon channel $g\to g$'', equation after label \texttt{nlogvquark}, quoted): $\mathcal{P}_{gg}(\xi)=2\big[\frac{\xi}{(1-\xi)_{+}}+\frac{1-\xi}{\xi}+\xi(1-\xi)\big]+\big(\frac{11}{6}-\frac{2N_{f}T_{R}}{3N_{c}}\big)\delta(1-\xi)$, entering as $N_{c}\mathcal{P}_{gg}$ in their Eqs.~(\texttt{gluonr}), (\texttt{gluonfragr}). For $N_{f}=0$: $N_{c}\mathcal{P}_{gg}=2C_{A}\big[\frac{\xi}{(1-\xi)_{+}}+\frac{1-\xi}{\xi}+\xi(1-\xi)\big]+\frac{\beta_{0}}{2}\delta(1-\xi)$, $\beta_{0}=\frac{11N_{c}}{3}$: the standard LO $P_{gg}$.

\noindent\textbf{Comparison.}
\begin{itemize}
\item \emph{Normalisation and $\delta(1-\xi)$ coefficient:} identical ($2N_{c}$ in front of the bracket, $\frac{\beta_{0}}{2}$ in front of $\delta(1-\xi)$). The value $\frac{\beta_{0}}{2}$ is the one required by the momentum sum rule $\int_{0}^{1}d\xi\,\xi P_{gg}(\xi)=0$ (verified numerically: $0$ to 20 digits). \textbf{AGREES.}
\item \emph{End point $\xi\to1$} (integrated gluon soft, $p^{+}\to\Lambda$): the prescription $\int d\xi\,\frac{\xi F(\xi)-F(1)}{1-\xi}$ is the standard $\xi\,[\frac{1}{1-\xi}]_{+}$. The subtracted piece is $\log\frac{k^{+}}{\Lambda}\,\mathcal{F}(1)$, the JIMWLK (target) log. CXY do the same: the rapidity divergence at $\xi\to1$ is subtracted into the BK/JIMWLK evolution of the target (CXY Eq.~(\texttt{BKg2})), after which ``we can replace $\frac{1}{1-\xi}$ by $\frac{1}{(1-\xi)_{+}}$''. \textbf{AGREES IN STRUCTURE.}
\item \emph{End point $\xi\to0$} (measured gluon soft relative to its parent, $p^{+}\to\vee$): the standard $P_{gg}$ has no subtraction there; $\widetilde{\mathcal P}_{gg}$ has $[\frac{1-\xi}{\xi}]_{+}F\equiv\frac{1-\xi}{\xi}[F(\xi)-F(0)]$. On $[\frac{k^{+}}{\vee},1]$ the subtracted piece is $-2N_{c}\mathcal{C}\,\mathcal{F}(0)\big[\log\frac{\vee}{k^{+}}-1+\frac{k^{+}}{\vee}\big]$, i.e.\ the BFKL term of the row, $+2N_{c}\mathcal{C}\frac{4\pi}{k^{2}}\log\frac{\vee}{k^{+}}\,\mathcal{X}$, and the remnant $-2N_{c}\mathcal{C}\frac{4\pi}{k^{2}}\mathcal{X}=-\frac{1}{(2\pi)^{3}}\frac{g^{4}N_{c}}{2\pi^{3}}\frac{1}{k^{2}k^{+}}\mathcal{X}$, which is the second term of Group~II Row~IV in the finite terms. Hence
\begin{equation}
N_{\rm DGLAP}+N_{\rm BFKL}^{\rm (II\,IV)}+N_{\rm remnant}=-\mathcal{C}\int_{k^{+}/\vee}^{1}d\xi\,P_{gg}^{\rm std}(\xi)\,\mathcal{F}(\xi)+{\cal O}\Big(\frac{k^{+}}{\vee}\Big)
\end{equation}
(verified with a model $\mathcal{F}$: difference $\to0$ like $k^{+}/\vee$; script [1d]). In CXY the projectile parton is collinear and $\xi\ge\tau/z$ is bounded from below, so no such subtraction arises; here the small-$\xi$ region is projectile (BFKL) evolution.
\item \emph{KLSZ} (2308.15545, Eqs.~(\texttt{sfat}), (\texttt{xi++}), (\texttt{integral})) use the ``double plus'' prescription $\int_{0}^{1}\frac{d\xi}{\xi_{+}(1-\xi)_{+}}N(\xi)f\equiv\int_{0}^{1}d\xi\big[\frac{N f}{\xi(1-\xi)}-\frac{f(0)}{\xi}-\frac{f(1)}{1-\xi}\big]$, $N=\xi^{2}+(1-\xi)^{2}+\xi^{2}(1-\xi)^{2}$, ``in contradistinction to the usual DGLAP splitting function'', because ``both poles are already included in the energy evolution generated by the LO JIMWLK equation''; $\int\frac{d\xi}{\xi_{+}(1-\xi)_{+}}N=-\frac{11}{6}$. Since $\frac{N}{\xi(1-\xi)}=\xi(1-\xi)+\frac{\xi}{1-\xi}+\frac{1-\xi}{\xi}$ and $N(0)=N(1)=1$:
\begin{equation}
\widetilde{\mathcal P}_{gg}\otimes F-2N_{c}F(0)=2N_{c}\int_{0}^{1}\frac{d\xi}{\xi_{+}(1-\xi)_{+}}N(\xi)F(\xi)+\frac{\beta_{0}}{2}F(1),
\end{equation}
i.e.\ TARGET's DGLAP term plus the remnant ($-\mathcal{C}[\widetilde{\mathcal P}_{gg}\otimes\mathcal{F}-2N_{c}\mathcal{F}(0)]$) is exactly the KLSZ double-plus function times $2N_{c}$, plus $\frac{\beta_{0}}{2}\delta(1-\xi)$ (numerically equal to 20 digits for a model $F$, [1a]). The reason for the double subtraction is the same as in KLSZ (here: BFKL at $\xi\to0$, JIMWLK at $\xi\to1$). \textbf{AGREES.}
\item \emph{$\delta(1-\xi)$ term vs.\ RC:} $-\mathcal{C}\frac{\beta_{0}}{2}\mathcal{F}(1)=-\frac{1}{(2\pi)^{3}}\frac{g^{2}}{2\pi^{2}}\frac{g^{2}\beta_{0}}{16\pi^{2}}\frac{1}{k^{+}}\int e^{-ik\cdot(y-y')}\big[2\gamma_{E}+\log\frac{(y'-y)^{2}\mu^{2}}{4}\big]\dots=-N_{\rm RC}$, identically (same prefactor, same bracket, same integrand). $\beta_{0}$ is \emph{not} a derived coefficient here: the real emission produces no $\delta(1-\xi)$, and the virtual $\frac{11}{3}$ terms cancel among themselves (Group~I Row~II, $+\frac{11}{3}$, against Row~V, $2\times(-\frac{11}{6})$). It is an add-and-subtract choice of presentation; its value $\frac{\beta_{0}}{2}$ is the standard one (momentum sum rule; consistency with $\overline{\rm MS}$ coupling and fragmentation function, below).
\end{itemize}

\subsection{The collinear logarithm: final-state $g\to g$ fragmentation}

\noindent\textbf{Form.} With $d^{3}q=\xi^{-3}d^{3}k$ ($q^{+}=k^{+}/\xi$, $\boldsymbol q=\boldsymbol k/\xi$) and $\frac{1}{q^{+}}=\frac{\xi}{k^{+}}$, $\frac{1}{k^{+}}\frac{1}{\xi^{2}}e^{-ik\cdot(y-y')/\xi}=\xi^{-3}\times\big[\frac{1}{q^{+}}e^{-iq\cdot(y-y')}\big]$, and $\mathcal{C}=\big[\frac{1}{(2\pi)^{3}}\frac{g^{2}}{2\pi^{2}}\frac{1}{k^{+}}\big]\times\frac{\alpha_{s}}{2\pi}$ ($\frac{g^{2}}{8\pi^{2}}=\frac{\alpha_{s}}{2\pi}$). So
\begin{equation}
\frac{d^{3}N_{\rm DGLAP}}{d^{3}k}=-\frac{\alpha_{s}}{2\pi}\int_{0}^{1}\frac{d\xi}{\xi^{3}}\,\widetilde{\mathcal P}_{gg}(\xi)\,\Big[\frac{d^{3}N_{\rm LO}}{d^{3}q}\ \text{with}\ \log\frac{(y'-y)^{2}\mu^{2}}{c_{0}^{2}\,\xi^{2}}\ \text{inserted}\Big]_{q=k/\xi},\qquad c_{0}=2e^{-\gamma_{E}},
\end{equation}
since $2\gamma_{E}+\log\frac{r^{2}\mu^{2}}{4}=\log\frac{r^{2}\mu^{2}}{c_{0}^{2}}$. This is the LO cross section of a parent gluon at $(k^{+}/\xi,k/\xi)$ convoluted with $P_{gg}$ and a collinear log: final-state fragmentation $g\to g$. Factor $\xi^{-3}$ ($=\frac{1}{\xi^{2}}$ with the $\frac{1}{k^{+}}$): correct; coefficient $\frac{\alpha_{s}}{2\pi}$ times the LO prefactor: correct; sign: the explicit $\mu$-dependence is $-\frac{\alpha_{s}}{2\pi}\log\mu^{2}\,P\otimes{\rm LO}$, which is what $D_{g/g}(\xi,\mu)=\delta(1-\xi)+\frac{\alpha_{s}}{2\pi}P_{gg}(\xi)\log\frac{\mu^{2}}{\mu_{0}^{2}}+\dots$ (standard, from memory; CXY Eq.~(\texttt{gluonfragr}) is the bare-to-renormalised form) compensates. \emph{Check of the $\mu$-dependence:} $\partial\mathcal{F}/\partial\log\mu^{2}=\mathcal{G}_{0}(\xi)$ and $\mathcal{G}_{0}(0)=0$, so the subtraction at $\xi=0$ does not touch it and $\partial N_{\rm DGLAP}/\partial\log\mu^{2}=-\mathcal{C}\int_{0}^{1}d\xi\,P^{\rm std}_{gg}\,\mathcal{G}_{0}$, the standard DGLAP kernel including its $\frac{\beta_{0}}{2}\delta(1-\xi)$; the $+\mathcal{C}\frac{\beta_{0}}{2}\mathcal{G}_{0}(1)\log\mu^{2}$ of $N_{\rm RC}$ compensates the running of $g^{2}(\mu)$ in the LO term. The result is therefore consistent with an $\overline{\rm MS}$ coupling and an $\overline{\rm MS}$ gluon fragmentation function at the same scale $\mu$. \textbf{AGREES.}

\noindent\textbf{CXY 1203.6139}, Eq.~after label \texttt{gluonfragr} (quoted):
\begin{equation*}
\mathcal{H}_{2gg}^{(1)}=N_{c}\mathcal{P}_{gg}(\xi)\ln\frac{c_{0}^{2}}{r_{\perp}^{2}\mu^{2}}\Big(e^{-ik_{\perp}\cdot r_{\perp}}+\frac{1}{\xi^{2}}e^{-i\frac{k_{\perp}}{\xi}\cdot r_{\perp}}\Big)-\Big(\frac{11}{3}-\frac{4N_{f}T_{R}}{3N_{c}}\Big)N_{c}\delta(1-\xi)e^{-ik_{\perp}\cdot r_{\perp}}\ln\frac{c_{0}^{2}}{r_{\perp}^{2}k_{\perp}^{2}},
\end{equation*}
multiplying $S^{(2)}(x,y)S^{(2)}(y,x)$ (large-$N_{c}$ adjoint dipole), with $r_{\perp}=x_{\perp}-y_{\perp}$, inside $\int\frac{dz}{z^{2}}\frac{dx}{x}\xi xg(x,\mu)D_{h/g}(z,\mu)$; collinear poles absorbed by Eqs.~(\texttt{gluonr}), (\texttt{gluonfragr}) with $c_{0}^{2}/(r^{2}\mu^{2})$; rapidity divergence at $\xi\to1$ into BK (Eq.~(\texttt{BKg2})). 1112.1061, Eq.~(\texttt{h2qq}), has the same structure in the $q\to q$ channel.
\begin{itemize}
\item \emph{Counterpart:} the term $\frac{1}{\xi^{2}}e^{-i\frac{k}{\xi}\cdot r}N_{c}\mathcal{P}_{gg}\ln\frac{c_{0}^{2}}{r^{2}\mu^{2}}$ (final state, CXY Fig.~\texttt{nlogr}(b), $S^{(2)}(v,v')S^{(2)}(v',v)$ with $v=\xi x+(1-\xi)b$: the same parent coordinate as TARGET's $y=z+\xi(w-z)$) corresponds to TARGET's $N_{\rm DGLAP}$: same $\frac{1}{\xi^{2}}$, sign, coefficient, splitting function (up to the $\xi\to0$ subtraction above) and the same $c_{0}$. \textbf{AGREES IN STRUCTURE.}
\item \emph{No counterpart:} (a) $e^{-ik\cdot r}N_{c}\mathcal{P}_{gg}\ln\frac{c_{0}^{2}}{r^{2}\mu^{2}}$ (initial-state DGLAP of the projectile PDF): the projectile here is $\rho$ with an unintegrated (transverse-momentum dependent) distribution; its small-$x$ evolution is the BFKL term, and there is no collinear pole of the initial state. (b) The virtual term $-\frac{11}{3}N_{c}\delta(1-\xi)\ln\frac{c_{0}^{2}}{r^{2}k^{2}}$: in this calculation the virtual $\frac{11}{3}$ terms cancel (Group~I Rows~II, V). (c) Large $N_{c}$ and sextupoles: Row~IV needs no large-$N_{c}$ limit here (the colour factor is exactly $N_{c}\times$LO structure).
\item \emph{Argument of the log:} TARGET's log is $\log\frac{u^{2}\mu^{2}}{c_{0}^{2}}$ with $u=\frac{y'-y}{\xi}$ the separation of the \emph{measured} gluon in amplitude and conjugate amplitude (conjugate to $k$); CXY's (and AABKL's $\overline{\rm MS}$ form, App.~C, Eq.~(\texttt{sigma\_coll\_MSbar\_bis})) is $\log\frac{r^{2}\mu^{2}}{c_{0}^{2}}$ with $r=y'-y$ the \emph{parent} separation (conjugate to $k/\xi$). Hence
\begin{equation}
N_{\rm DGLAP}=N_{\rm DGLAP}^{\rm CXY\text{-}form}+\mathcal{C}\int_{0}^{1}d\xi\,\widetilde{\mathcal P}_{gg}(\xi)\,\log\xi^{2}\;\mathcal{G}_{0}(\xi),
\end{equation}
a finite term (zero at $\xi=1$, so the $\delta$ term and $N_{\rm RC}$ are unaffected; convergent at $\xi\to0$; $\mathcal{F}$ and the CXY-form function have the same value at $\xi=0$, checked numerically, [2b], [2c]). It is a convention of the $d$-dimensional convolution (derived here): in TARGET only the daughter separation $w-z$ at fixed $u$ is continued to $d=2-\epsilon$, so the pole term is $\frac{2}{\epsilon}\xi^{-2}\times$(2-dimensional LO at $k/\xi$) and the finite log is that of $u$. If $u$ is also continued, $d^{d}u\,(u^{2})^{\epsilon/2}=\xi^{-2}d^{d}r\,(r^{2})^{\epsilon/2}$; a counterterm $\frac{2}{\epsilon}\xi^{-2}{\rm LO}_{d}(k/\xi)$ then leaves $\log\frac{r^{2}\mu^{2}}{c_{0}^{2}}$ (CXY form), a counterterm with the $d$-dimensional Jacobian $\frac{2}{\epsilon}\xi^{-d}{\rm LO}_{d}(k/\xi)$ leaves $\log\frac{u^{2}\mu^{2}}{c_{0}^{2}}$ (TARGET form). Which one is the $\overline{\rm MS}$ fragmentation function of the fits depends on the $d$-dimensional definition of $D_{h/g}$ (AABKL App.~C use $\zeta^{-(D-2)}$ and $(1-\xi)^{-(5-D)}$ and obtain the parent form); this was not re-derived here. AABKL's own cut-off scheme (main text, Eqs.~(\texttt{fact1}), (\texttt{g2ggRF}): cut $|z|>1/\mu$ on the daughter separation at fixed measured separation) has the log of the measured separation, as TARGET: $\int d^{2}z\big[\frac{z\cdot(z+u)}{z^{2}(z+u)^{2}}-\frac{\theta(|z|\mu_{\rm CO}-1)}{z^{2}}\big]=-\pi\log(u^{2}\mu_{\rm CO}^{2})$ (numerically $-1.931983347638915$ against $-1.931983347638913$ for $u=0.8$, $\mu_{\rm CO}=1.7$, [2a]), so TARGET's bracket equals AABKL's cut-off log with $\mu_{\rm CO}=\mu/c_{0}$, i.e.\ $\mu^{2}_{\overline{\rm MS}}=c_{0}^{2}\mu_{\rm CO}^{2}=1.26\,\mu^{2}_{\rm CO}$, exactly AABKL's conversion ``$f_{\mu^{2}}=f^{\overline{MS}}_{1.26\mu^{2}}$'' (App.~C, Eqs.~(\texttt{C\_X\_result}), (\texttt{subtr\_MSbar})). \textbf{AGREES IN STRUCTURE; finite $\widetilde{\mathcal P}_{gg}\log\xi^{2}\otimes$LO term is scheme dependent.} To use fitted $\overline{\rm MS}$ fragmentation functions, either move this term to the finite terms or state the convention.
\item \emph{The pole.} The collinear pole $\frac{2}{\epsilon}$ of Row~IV is not written (rule $\int d^{2}p/p^{2}=0$). In CXY and AABKL it is absorbed into $D_{h/g}$. KW 0712.3732 (abstract): the gluon production cross section ``is not infrared-safe''; for hadrons the collinear divergence is absorbed into a fragmentation function. TARGET's result is therefore a partonic coefficient in a definite collinear scheme; it should be convoluted with $D_{h/g}(z,\mu)$ (or the measured gluon should be defined at resolution $\mu$).
\end{itemize}

\subsection{Running coupling}

\noindent\textbf{Scale.} One-loop running (HK Eq.~(\texttt{as\_geom}) with $\beta_{2}=\beta_{0}/4\pi$, KLSZ Eq.~(\texttt{beta})): $\alpha_{s}(Q^{2})=\alpha_{s}(\mu^{2})/[1+\frac{\alpha_{s}\beta_{0}}{4\pi}\log\frac{Q^{2}}{\mu^{2}}]$. The LO+RC bracket $1+\frac{\alpha_{s}\beta_{0}}{4\pi}\log\frac{r^{2}\mu^{2}e^{2\gamma_{E}}}{4}$ is $\alpha_{s}(Q^{2})/\alpha_{s}(\mu^{2})$ to this order with $Q^{2}=\frac{4e^{-2\gamma_{E}}}{r^{2}}=\frac{c_{0}^{2}}{r^{2}}$, $r=y'-y$; sign: asymptotically free. But because $N_{\rm RC}=-(\delta\text{-term of }N_{\rm DGLAP})$,
\begin{equation}
N_{\rm LO}+N_{\rm RC}+N_{\rm DGLAP}\big|_{\delta(1-\xi)}=g^{2}(\mu)\times[\text{LO}],
\end{equation}
so the scale $c_{0}^{2}/r^{2}$ is the choice $\mu=c_{0}/r$ in the $\delta(1-\xi)$ part of the collinear log, the same choice CXY make for the factorisation scale (1203.6139 after Eq.~(\texttt{hgg6}); 1112.1061 after Eq.~(\texttt{h2qq})). Any other argument in the subtracted $\delta$ term is equally allowed and shifts a finite $\frac{\alpha_{s}\beta_{0}}{4\pi}\log(\dots)\times$LO between the two terms.

\noindent\textbf{HK 1009.0545, KW 0712.3732.} KW Eq.~(\texttt{Gmult3}): $\frac{dN^{G}}{d^{2}q\,dy}=\frac{1}{(2\pi)^{2}}|\boldsymbol M^{a}(\boldsymbol q,q_{+})|^{2}\alpha_{s}(\Lambda_{\rm coll}^{2})$: ``the bare coupling constant stemming from the vertices of gluon emission by the projectile in the amplitude and in the complex conjugate amplitude gets renormalized'' at the collinear resolution scale $\Lambda_{\rm coll}$. HK Eq.~(\texttt{brem\_rc}): $\alpha_{\mu}^{3}\to\alpha_{s}(\Lambda^{2}_{\rm coll}e^{-5/3})\,\alpha_{s}^{2}(\boldsymbol q^{2}e^{-5/3})$, final result Eq.~(\texttt{rc\_incl}): $\frac{d\sigma}{d^{2}k\,dy}=\frac{2C_{F}}{\pi^{2}}\frac{\alpha_{s}(\Lambda_{\rm coll}^{2}e^{-5/3})}{\boldsymbol k^{2}}\int\frac{d^{2}q}{\boldsymbol q^{2}(\boldsymbol k-\boldsymbol q)^{2}}\frac{\alpha_{s}^{2}(\boldsymbol q^{2}e^{-5/3})\alpha_{s}^{2}((\boldsymbol k-\boldsymbol q)^{2}e^{-5/3})}{\alpha_{s}(Q^{2}e^{-5/3})\alpha_{s}(Q^{*2}e^{-5/3})}$; HK state that $\Lambda_{\rm coll}$ ``plays the role of the factorization scale for the fragmentation functions''.
\begin{itemize}
\item \emph{Comparable:} the coupling of the emission vertex of the measured gluon. In TARGET it is $g^{2}(\mu)$ with $\mu$ the scale of the collinear (fragmentation) subtraction, once $N_{\rm RC}$ is regrouped with the $\delta$ term; KW/HK: $\alpha_{s}(\Lambda^{2}_{\rm coll})$. Moreover the $\frac{11}{3}\log[(x-w)^{2}\mu^{2}]$ terms of Group~I (emitter--gluon distance) cancel between Rows~II and V, i.e.\ the emission coupling does not run with the emitter--gluon distance; this is what KW find for the produced gluon (the $\boldsymbol q$ dependence of the bubble chain on the cut line disappears, Eqs.~(\texttt{Gmult-1})--(\texttt{Gmult3})). \textbf{AGREES IN STRUCTURE.} The constants are not comparable ($e^{-5/3}$ is the $\overline{\rm MS}$ quark-bubble constant of the BLM procedure; TARGET has $N_{f}=0$ and no bubble).
\item \emph{Not comparable:} HK's other four couplings and the complex scales $Q,Q^{*}$ (Eq.~(\texttt{Qscale})) belong to the $t$-channel exchanges, which here are inside $\rho\rho$ and $U$ and are not renormalised in this calculation; $\beta_{0}$ enters TARGET once, at the emission vertex, and only through the $\delta(1-\xi)$ term. KW hep-ph/0612071 (Eqs.~(\texttt{Kall}), (\texttt{Q})) gives the triumvirate $\alpha_{s}(\boldsymbol q^{2}e^{-5/3})\alpha_{s}(\boldsymbol q'^{2}e^{-5/3})/\alpha_{s}(Q^{2}e^{-5/3})$ for the (virtual) gluon of one step of small-$x$ evolution; it applies to the BFKL/JIMWLK kernels, not to the measured gluon.
\item \emph{Dilute limit, momentum space.} With $\log\frac{r^{2}\mu^{2}}{c_{0}^{2}}=-\frac{1}{\pi}\int d^{2}q\,\frac{e^{-iq\cdot r}-\theta(\mu-q)}{q^{2}}$ (verified), the RC term in momentum space is
\begin{equation}
N_{\rm RC}(k)=-\frac{\alpha_{s}\beta_{0}}{4\pi}\,\frac{1}{\pi}\int d^{2}q\,\frac{L(\boldsymbol k+\boldsymbol q)-\theta(\mu-q)L(\boldsymbol k)}{q^{2}},
\end{equation}
$L(\boldsymbol k)$ the LO spectrum (prefactor and $\int_{Y}$ included). It is a convolution, not $L(k)$ times a coupling at a local scale. For a power law $L\propto k^{-2-2a}$ one finds $Q_{\rm eff}^{2}=k^{2}e^{-[\psi(1+a)+\psi(-a)+2\gamma_{E}]}$ (e.g.\ $64k^{2}$ for $a=-\frac14$; the momentum-space integral gives $64.30$, $0.5\%$ quadrature error, [3a]). For a spectrum falling faster than $k^{-2}$ (dilute, $L\sim k^{-4}$) the region $\boldsymbol k+\boldsymbol q$ soft dominates: $N_{\rm RC}(k)\to-\frac{\alpha_{s}\beta_{0}}{4\pi}\frac{1}{\pi k^{2}}\int d^{2}p\,L(p)$, a $1/k^{2}$ tail that dominates $L\sim1/k^{4}$, so no local scale exists (model $L=(k^{2}+m^{2})^{-2}$: $N_{\rm RC}(k)$ divided by this tail $=1.13,\,1.06,\,1.012$ for $k/m=3,10,30$, and $\langle\log\frac{r^{2}m^{2}}{c_{0}^{2}}\rangle=-12.6,\,-108,\,-913$, i.e.\ a would-be $Q_{\rm eff}^{2}=\mu^{2}e^{-\langle\log\rangle}$ grows like $e^{k^{2}/m^{2}}$; model $L=\frac{\log(1+k^{2}/m^{2})}{k^{2}(k^{2}+m^{2})}$: ratios $1.08,\,1.12,\,1.04$; [3b]). This tail is $\frac{\beta_{0}}{2}\mathcal{C}\mathcal{F}(0)$, the BFKL structure, and cancels against the $\delta$ term of $N_{\rm DGLAP}$. Hence HK's local scales cannot be compared with $c_{0}^{2}/r^{2}$ term by term; only the regrouped statement above is comparable.
\end{itemize}

\noindent\textbf{KLSZ 2308.15545.} Their statement (Sec.~``The dressed gluon and the DGLAP splittings''): in NLO JIMWLK the coefficient of $\ln\mu^{2}_{\overline{MS}}$ in $K_{JSJ}$ is $2\beta_{0}$ relative to LO; one $\beta_{0}$ belongs to the running coupling (via ${\cal Z}^{1/2}(Q^{2})=1+\frac{\alpha_{s}}{8\pi}\beta_{0}\ln\frac{Q^{2}}{\mu^{2}}$, Eq.~(\texttt{Z}), with the natural scale the emitter distance $X$, Eq.~(\texttt{alpharesummed})), the other to the DGLAP-like dressing of the gluon (Eqs.~(\texttt{sfat1}), (\texttt{subst})); ``the UV logarithm is associated with setting the scale of the running coupling, while the $Q^{2}$ dependent logarithm is associated with setting the resolution scale for DGLAP splittings''. In TARGET the UV $\beta_{0}$ logs with the emitter distance cancel (Group~I), and the only $\beta_{0}$ is the $\delta(1-\xi)$ of the final-state $P_{gg}$, written once as $-$ (DGLAP) and once as $+$ (RC). By KLSZ's classification TARGET's $\beta_{0}$ terms are DGLAP (resolution) logs, and the RC term is not a derived running of the coupling. Formal remark: if $\mathcal{F}(\xi)$ is replaced by $\mathcal{F}(1)$ (unresolved final state), $N_{\rm DGLAP}+N_{\rm remnant}\to+\mathcal{C}\frac{\beta_{0}}{2}\mathcal{F}(1)=N_{\rm RC}$, so the $\beta_{0}$ terms of the row are $2\times N_{\rm RC}$, the same ``$2\beta_{0}$'' and the same split (one $\beta_{0}$ called RC, one DGLAP); the same counting gives the coefficients $-\frac{11}{12}$, $-\frac{11}{12}$ of the large-distance tail in \texttt{Finite\_terms\_final} (OPEN II/III-1). \emph{Verdict on ``RC comes as a result of DGLAP'':} correct as a statement about this calculation (RC $=-\delta$ term of DGLAP, added and subtracted); consistent with KLSZ if it means ``the $\beta_{0}$ here has collinear (DGLAP) origin''; not consistent with KLSZ if it means that DGLAP generates the running of the coupling: in KLSZ running coupling and DGLAP are distinct, and the running-coupling scale is set by the UV log at the emitter distance, which is absent here. Suggested wording: ``Writing the collinear term with the standard $P_{gg}$, whose $\delta(1-\xi)$ coefficient is $\beta_{0}/2$, requires the compensating term $N_{\rm RC}$; with it, the LO coupling may be read as $g^{2}(c_{0}^{2}/r^{2})$, equivalently as $g^{2}(\mu)$ at the fragmentation scale (cf.\ KW, HK: $\alpha_{s}(\Lambda^{2}_{\rm coll})$).''

\subsection{Summary of this section}
\begin{center}
\begin{tabular}{|p{5.2cm}|p{4.6cm}|p{7.2cm}|}\hline
Item & Source & Outcome\\\hline
Normalisation, $\delta$ coefficient $\frac{\beta_{0}}{2}$ & CXY 1203.6139 & AGREES\\
$\xi\to1$ plus prescription & CXY (BK subtraction) & AGREES IN STRUCTURE (JIMWLK log)\\
$\xi\to0$ subtraction & KLSZ double plus & AGREES (identity with the remnant); absent in CXY\\
$\delta$ term $=-N_{\rm RC}$ & derived here & exact; $\beta_{0}$ added and subtracted\\
$\frac{1}{\xi^{2}}$, sign, $\frac{\alpha_{s}}{2\pi}$, $c_{0}$, $\mu$-dependence & CXY, AABKL; derived & AGREES\\
log argument & CXY, AABKL App.~C & differs: measured vs parent separation; finite $\widetilde{\mathcal P}\log\xi^{2}\otimes$LO (scheme); equals AABKL cut-off scheme with $\mu_{\rm CO}=\mu/c_{0}$\\
initial-state DGLAP, virtual $\frac{11}{3}\ln\frac{c_{0}^{2}}{r^{2}k^{2}}$ & CXY & no counterpart (explained)\\
scale of emission coupling & KW 0712.3732, HK & AGREES IN STRUCTURE ($g^{2}(\mu_{\rm frag})\leftrightarrow\alpha_{s}(\Lambda^{2}_{\rm coll})$); $c_{0}^{2}/r^{2}$ is a choice\\
HK $t$-channel couplings, KW triumvirate & HK, KW & COULD NOT BE COMPARED (not in this calculation)\\
``RC from DGLAP'' & KLSZ & consistent only as ``$\beta_{0}$ of collinear origin''\\\hline
\end{tabular}
\end{center}
